\documentclass[12pt,a4paper]{article}
\usepackage[UTF8]{ctex}
\usepackage{amsmath,amssymb,amsthm}
\usepackage{geometry}
\usepackage{bm}

\geometry{margin=2.5cm}

\title{欧拉方程证明与拓展：$m_b = I_b \dot{\omega}_b + [\omega_b] I_b \omega_b$}
\author{第8章补充说明}
\date{}

\begin{document}

\maketitle

\section{问题提出}

\textbf{问题}：证明并拓展欧拉方程：
\begin{equation}
m_b = I_b \dot{\omega}_b + [\omega_b] I_b \omega_b
\label{eq:euler_equation}
\end{equation}

\textbf{符号说明}：
\begin{itemize}
\item $m_b \in \mathbb{R}^3$：作用在刚体上的力矩（在体坐标系$\{b\}$中表示）
\item $I_b \in \mathbb{R}^{3 \times 3}$：转动惯量矩阵（在体坐标系$\{b\}$中）
\item $\omega_b \in \mathbb{R}^3$：角速度（在体坐标系$\{b\}$中表示）
\item $\dot{\omega}_b \in \mathbb{R}^3$：角加速度（在体坐标系$\{b\}$中表示）
\item $[\omega_b]$：角速度$\omega_b$的反对称矩阵
\end{itemize}

\section{欧拉方程的历史背景}

\subsection{欧拉（Leonhard Euler）}

欧拉方程由瑞士数学家欧拉（1707-1783）提出，用于描述刚体的旋转动力学。

\subsection{方程的重要性}

欧拉方程是刚体动力学的核心方程之一，与牛顿方程一起构成了刚体动力学的完整描述。

\section{证明方法1：从角动量定理出发}

\subsection{角动量的定义}

刚体的角动量（在体坐标系$\{b\}$中）定义为：
\begin{equation}
\bm{L}_b = I_b \omega_b
\label{eq:angular_momentum}
\end{equation}

\textbf{展开形式}：
\begin{equation}
\bm{L}_b = \begin{bmatrix}
I_{xx} & I_{xy} & I_{xz} \\
I_{xy} & I_{yy} & I_{yz} \\
I_{xz} & I_{yz} & I_{zz}
\end{bmatrix} \begin{bmatrix}
\omega_x \\
\omega_y \\
\omega_z
\end{bmatrix}
\end{equation}

\subsection{角动量定理}

根据角动量定理，力矩等于角动量的时间导数：
\begin{equation}
m_b = \frac{d\bm{L}_b}{dt}
\label{eq:angular_momentum_theorem}
\end{equation}

\textbf{关键点}：这里的导数是在\textbf{惯性坐标系}中计算的，而$\bm{L}_b = I_b \omega_b$是在\textbf{体坐标系}中表示的。

\subsection{旋转坐标系中的时间导数}

对于在旋转坐标系中表示的向量$\bm{L}_b$，其在惯性系中的时间导数为：
\begin{equation}
\left.\frac{d\bm{L}_b}{dt}\right|_{\text{inertial}} = \left.\frac{d\bm{L}_b}{dt}\right|_{\text{body}} + \omega_b \times \bm{L}_b
\label{eq:rotating_frame_derivative}
\end{equation}

\textbf{物理意义}：
\begin{itemize}
\item 第一项：在体坐标系中观察到的角动量变化率
\item 第二项：由于坐标系旋转产生的额外项（科里奥利效应）
\end{itemize}

\subsection{完成证明}

将式(\ref{eq:angular_momentum})和式(\ref{eq:rotating_frame_derivative})代入式(\ref{eq:angular_momentum_theorem})：
\begin{align}
m_b &= \frac{d}{dt}(I_b \omega_b) + \omega_b \times (I_b \omega_b) \\
&= I_b \frac{d\omega_b}{dt} + \frac{dI_b}{dt} \omega_b + \omega_b \times (I_b \omega_b)
\end{align}

\textbf{关键观察}：由于$I_b$在体坐标系中是常数（刚体的形状和质量分布不变），$\frac{dI_b}{dt} = 0$。

\textbf{因此}：
\begin{equation}
m_b = I_b \dot{\omega}_b + \omega_b \times (I_b \omega_b)
\end{equation}

\textbf{使用反对称矩阵}：
\begin{equation}
\omega_b \times (I_b \omega_b) = [\omega_b] I_b \omega_b
\end{equation}

\textbf{因此}：
\begin{equation}
m_b = I_b \dot{\omega}_b + [\omega_b] I_b \omega_b
\end{equation}

\textbf{证毕}。

\section{证明方法2：从力矩的定义出发}

\subsection{力矩的定义}

对于由点质量组成的刚体，力矩定义为：
\begin{equation}
m_b = \sum_i [r_i] m_i \ddot{r}_i
\label{eq:torque_definition}
\end{equation}

其中：
\begin{itemize}
\item $r_i$：第$i$个点质量在体坐标系$\{b\}$中的位置
\item $m_i$：第$i$个点质量的质量
\item $\ddot{r}_i$：第$i$个点质量的加速度
\end{itemize}

\subsection{点质量的加速度}

对于刚体上的点质量，其加速度为：
\begin{equation}
\ddot{r}_i = \dot{v}_b + [\dot{\omega}_b] r_i + [\omega_b] v_b + [\omega_b]^2 r_i
\label{eq:point_acceleration}
\end{equation}

其中：
\begin{itemize}
\item $\dot{v}_b$：质心的线加速度
\item $[\dot{\omega}_b] r_i$：由于角加速度产生的切向加速度
\item $[\omega_b] v_b$：由于线速度和角速度耦合产生的加速度
\item $[\omega_b]^2 r_i$：由于旋转产生的向心加速度
\end{itemize}

\subsection{展开力矩}

将式(\ref{eq:point_acceleration})代入式(\ref{eq:torque_definition})：
\begin{align}
m_b &= \sum_i [r_i] m_i (\dot{v}_b + [\dot{\omega}_b] r_i + [\omega_b] v_b + [\omega_b]^2 r_i) \\
&= \sum_i m_i [r_i] \dot{v}_b + \sum_i m_i [r_i][\dot{\omega}_b] r_i + \sum_i m_i [r_i][\omega_b] v_b + \sum_i m_i [r_i][\omega_b]^2 r_i
\end{align}

\subsection{简化各项}

\textbf{第一项}：$\sum_i m_i [r_i] \dot{v}_b$

如果质心在原点，$\sum_i m_i r_i = 0$，因此：
\begin{equation}
\sum_i m_i [r_i] \dot{v}_b = [\sum_i m_i r_i] \dot{v}_b = 0
\end{equation}

\textbf{第二项}：$\sum_i m_i [r_i][\dot{\omega}_b] r_i$

使用恒等式$[r_i][\dot{\omega}_b] r_i = -[r_i]^2 \dot{\omega}_b$：
\begin{equation}
\sum_i m_i [r_i][\dot{\omega}_b] r_i = -\sum_i m_i [r_i]^2 \dot{\omega}_b = I_b \dot{\omega}_b
\end{equation}

其中$I_b = -\sum_i m_i [r_i]^2$是转动惯量矩阵。

\textbf{第三项}：$\sum_i m_i [r_i][\omega_b] v_b$

如果质心在原点，这一项为零。

\textbf{第四项}：$\sum_i m_i [r_i][\omega_b]^2 r_i$

使用恒等式$[r_i][\omega_b]^2 r_i = -[\omega_b][r_i]^2 \omega_b$：
\begin{equation}
\sum_i m_i [r_i][\omega_b]^2 r_i = -[\omega_b] \sum_i m_i [r_i]^2 \omega_b = [\omega_b] I_b \omega_b
\end{equation}

\subsection{合并结果}

\begin{align}
m_b &= 0 + I_b \dot{\omega}_b + 0 + [\omega_b] I_b \omega_b \\
&= I_b \dot{\omega}_b + [\omega_b] I_b \omega_b
\end{align}

\textbf{证毕}。

\section{各项的物理意义}

\subsection{第一项：$I_b \dot{\omega}_b$}

\textbf{名称}：惯性力矩（inertial torque）

\textbf{物理意义}：
\begin{itemize}
\item 由于角加速度产生的力矩
\item 与角加速度成正比
\item 类似于线运动中的$m \dot{v}$（惯性力）
\end{itemize}

\textbf{展开形式}：
\begin{equation}
I_b \dot{\omega}_b = \begin{bmatrix}
I_{xx} \dot{\omega}_x + I_{xy} \dot{\omega}_y + I_{xz} \dot{\omega}_z \\
I_{xy} \dot{\omega}_x + I_{yy} \dot{\omega}_y + I_{yz} \dot{\omega}_z \\
I_{xz} \dot{\omega}_x + I_{yz} \dot{\omega}_y + I_{zz} \dot{\omega}_z
\end{bmatrix}
\end{equation}

\subsection{第二项：$[\omega_b] I_b \omega_b$}

\textbf{名称}：陀螺力矩（gyroscopic torque）或科里奥利力矩（Coriolis torque）

\textbf{物理意义}：
\begin{itemize}
\item 由于旋转产生的力矩
\item 与角速度的平方成正比
\item 即使角速度恒定（$\dot{\omega}_b = 0$），只要$\omega_b \neq 0$，就需要力矩来维持旋转
\item 这是陀螺效应的体现
\end{itemize}

\textbf{展开形式}：
\begin{align}
[\omega_b] I_b \omega_b &= \begin{bmatrix}
0 & -\omega_z & \omega_y \\
\omega_z & 0 & -\omega_x \\
-\omega_y & \omega_x & 0
\end{bmatrix} \begin{bmatrix}
I_{xx} \omega_x + I_{xy} \omega_y + I_{xz} \omega_z \\
I_{xy} \omega_x + I_{yy} \omega_y + I_{yz} \omega_z \\
I_{xz} \omega_x + I_{yz} \omega_y + I_{zz} \omega_z
\end{bmatrix} \\
&= \begin{bmatrix}
\omega_y (I_{xz} \omega_x + I_{yz} \omega_y + I_{zz} \omega_z) - \omega_z (I_{xy} \omega_x + I_{yy} \omega_y + I_{yz} \omega_z) \\
\omega_z (I_{xx} \omega_x + I_{xy} \omega_y + I_{xz} \omega_z) - \omega_x (I_{xz} \omega_x + I_{yz} \omega_y + I_{zz} \omega_z) \\
\omega_x (I_{xy} \omega_x + I_{yy} \omega_y + I_{yz} \omega_z) - \omega_y (I_{xx} \omega_x + I_{xy} \omega_y + I_{xz} \omega_z)
\end{bmatrix}
\end{align}

\section{主惯性轴情况}

\subsection{对角化形式}

如果坐标轴与主惯性轴对齐，转动惯量矩阵为对角矩阵：
\begin{equation}
I_b = \begin{bmatrix}
I_{xx} & 0 & 0 \\
0 & I_{yy} & 0 \\
0 & 0 & I_{zz}
\end{bmatrix}
\end{equation}

\subsection{欧拉方程的简化}

此时欧拉方程简化为：
\begin{equation}
\begin{bmatrix}
m_x \\
m_y \\
m_z
\end{bmatrix} = \begin{bmatrix}
I_{xx} \dot{\omega}_x \\
I_{yy} \dot{\omega}_y \\
I_{zz} \dot{\omega}_z
\end{bmatrix} + \begin{bmatrix}
(I_{zz} - I_{yy}) \omega_y \omega_z \\
(I_{xx} - I_{zz}) \omega_x \omega_z \\
(I_{yy} - I_{xx}) \omega_x \omega_y
\end{bmatrix}
\label{eq:euler_principal_axes}
\end{equation}

\textbf{展开形式}：
\begin{align}
m_x &= I_{xx} \dot{\omega}_x + (I_{zz} - I_{yy}) \omega_y \omega_z \\
m_y &= I_{yy} \dot{\omega}_y + (I_{xx} - I_{zz}) \omega_x \omega_z \\
m_z &= I_{zz} \dot{\omega}_z + (I_{yy} - I_{xx}) \omega_x \omega_y
\end{align}

\subsection{特殊情况}

\textbf{情况1}：绕$x$轴旋转（$\omega_y = \omega_z = 0$）
\begin{equation}
m_x = I_{xx} \dot{\omega}_x, \quad m_y = 0, \quad m_z = 0
\end{equation}

\textbf{情况2}：绕$y$轴旋转（$\omega_x = \omega_z = 0$）
\begin{equation}
m_x = 0, \quad m_y = I_{yy} \dot{\omega}_y, \quad m_z = 0
\end{equation}

\textbf{情况3}：绕$z$轴旋转（$\omega_x = \omega_y = 0$）
\begin{equation}
m_x = 0, \quad m_y = 0, \quad m_z = I_{zz} \dot{\omega}_z
\end{equation}

\textbf{结论}：绕主惯性轴旋转时，只需要一个方向的力矩，且陀螺项为零。

\section{能量关系}

\subsection{旋转动能}

刚体的旋转动能为：
\begin{equation}
K_{\text{rot}} = \frac{1}{2} \omega_b^T I_b \omega_b
\label{eq:rotational_kinetic_energy}
\end{equation}

\subsection{功率}

力矩的功率为：
\begin{equation}
P = m_b^T \omega_b
\end{equation}

\textbf{展开}：
\begin{align}
P &= (I_b \dot{\omega}_b + [\omega_b] I_b \omega_b)^T \omega_b \\
&= \dot{\omega}_b^T I_b^T \omega_b + ([\omega_b] I_b \omega_b)^T \omega_b \\
&= \dot{\omega}_b^T I_b \omega_b + (I_b \omega_b)^T [\omega_b]^T \omega_b \\
&= \dot{\omega}_b^T I_b \omega_b - (I_b \omega_b)^T [\omega_b] \omega_b \quad \text{（因为$[\omega_b]^T = -[\omega_b]$）} \\
&= \dot{\omega}_b^T I_b \omega_b - (I_b \omega_b)^T (\omega_b \times \omega_b) \\
&= \dot{\omega}_b^T I_b \omega_b - 0 \\
&= \dot{\omega}_b^T I_b \omega_b
\end{align}

\textbf{与动能的关系}：
\begin{equation}
P = \frac{d}{dt}\left(\frac{1}{2} \omega_b^T I_b \omega_b\right) = \frac{dK_{\text{rot}}}{dt}
\end{equation}

这验证了能量守恒。

\section{坐标变换}

\subsection{从体坐标系到空间坐标系}

如果力矩在空间坐标系中表示为$m_s$，角速度在空间坐标系中表示为$\omega_s$，那么：
\begin{equation}
m_s = R m_b, \quad \omega_s = R \omega_b
\end{equation}

其中$R$是从体坐标系到空间坐标系的旋转矩阵。

\textbf{空间坐标系中的欧拉方程}：
\begin{equation}
m_s = I_s \dot{\omega}_s + [\omega_s] I_s \omega_s
\end{equation}

其中$I_s = R I_b R^T$是空间坐标系中的转动惯量矩阵。

\section{应用示例}

\subsection{示例1：陀螺仪}

陀螺仪是一个快速旋转的刚体。即使没有外力矩，由于陀螺效应，陀螺仪会表现出进动。

\textbf{无外力矩情况}（$m_b = 0$）：
\begin{equation}
I_b \dot{\omega}_b + [\omega_b] I_b \omega_b = 0
\end{equation}

即使$\dot{\omega}_b = 0$（角速度恒定），只要$\omega_b \neq 0$且不沿主惯性轴，角速度方向会改变（进动）。

\subsection{示例2：航天器姿态控制}

在航天器姿态控制中，欧拉方程用于：
\begin{itemize}
\item 计算维持特定角速度所需的力矩
\item 设计姿态控制算法
\item 分析姿态稳定性
\end{itemize}

\subsection{示例3：机器人连杆动力学}

在机器人学中，每个连杆的旋转动力学由欧拉方程描述：
\begin{equation}
m_i = I_i \dot{\omega}_i + [\omega_i] I_i \omega_i
\end{equation}

这是递归牛顿-欧拉算法的基础。

\section{数值求解}

\subsection{给定力矩，求解角加速度}

如果已知力矩$m_b$和当前角速度$\omega_b$，可以求解角加速度：
\begin{equation}
\dot{\omega}_b = I_b^{-1} (m_b - [\omega_b] I_b \omega_b)
\end{equation}

\textbf{注意}：$I_b$是正定矩阵，因此可逆。

\subsection{给定角加速度，求解力矩}

如果已知角加速度$\dot{\omega}_b$和角速度$\omega_b$，可以直接计算力矩：
\begin{equation}
m_b = I_b \dot{\omega}_b + [\omega_b] I_b \omega_b
\end{equation}

\section{稳定性分析}

\subsection{无外力矩情况}

当$m_b = 0$时，欧拉方程为：
\begin{equation}
I_b \dot{\omega}_b = -[\omega_b] I_b \omega_b
\end{equation}

\textbf{绕主惯性轴旋转}：
\begin{itemize}
\item 如果$\omega_b$沿着主惯性轴方向，则$I_b \omega_b = \lambda \omega_b$（其中$\lambda$是主转动惯量）
\item 此时$[\omega_b] I_b \omega_b = [\omega_b] (\lambda \omega_b) = \lambda (\omega_b \times \omega_b) = 0$
\item 因此$I_b \dot{\omega}_b = 0$，即$\dot{\omega}_b = 0$（角速度恒定）
\item 这是稳定的旋转状态
\end{itemize}

\textbf{非主惯性轴旋转}：
\begin{itemize}
\item 如果$\omega_b$不沿着主惯性轴方向，则$[\omega_b] I_b \omega_b \neq 0$
\item 即使初始$\dot{\omega}_b = 0$，角速度方向会改变（进动）
\item 这是不稳定的旋转状态
\end{itemize}

\subsection{能量稳定性}

旋转动能$K_{\text{rot}} = \frac{1}{2} \omega_b^T I_b \omega_b$在无外力矩情况下守恒。

\textbf{绕主惯性轴旋转}：
\begin{itemize}
\item 动能恒定，角速度大小和方向都恒定
\item 这是能量最小的稳定状态
\end{itemize}

\textbf{非主惯性轴旋转}：
\begin{itemize}
\item 动能恒定，但角速度方向会改变
\item 角速度在空间中进动，但动能保持不变
\end{itemize}

\section{与牛顿方程的关系}

\subsection{完整的刚体动力学}

刚体的完整动力学由牛顿方程和欧拉方程共同描述：

\textbf{牛顿方程}（线运动）：
\begin{equation}
f_b = m \dot{v}_b + m [\omega_b] v_b
\label{eq:newton_equation}
\end{equation}

\textbf{欧拉方程}（旋转运动）：
\begin{equation}
m_b = I_b \dot{\omega}_b + [\omega_b] I_b \omega_b
\label{eq:euler_equation_complete}
\end{equation}

其中$m$是刚体的总质量。

\subsection{耦合关系}

线运动和旋转运动通过以下方式耦合：
\begin{itemize}
\item 线速度$v_b$和角速度$\omega_b$在牛顿方程中耦合
\item 线加速度$\dot{v}_b$和角加速度$\dot{\omega}_b$在力矩计算中耦合
\item 质心位置影响转动惯量矩阵（通过Steiner定理）
\end{itemize}

\section{矩阵形式}

\subsection{空间惯性矩阵}

刚体的完整动力学可以用$6 \times 6$空间惯性矩阵表示：
\begin{equation}
G_b = \begin{bmatrix}
I_b & 0 \\
0 & m I
\end{bmatrix}
\end{equation}

其中$I$是$3 \times 3$单位矩阵。

\subsection{统一的动力学方程}

使用空间速度$\mathcal{V}_b = \begin{bmatrix} \omega_b \\ v_b \end{bmatrix}$和空间力$\mathcal{F}_b = \begin{bmatrix} m_b \\ f_b \end{bmatrix}$：

\begin{equation}
\mathcal{F}_b = G_b \dot{\mathcal{V}}_b + [\mathcal{V}_b]^T G_b \mathcal{V}_b
\end{equation}

其中$[\mathcal{V}_b]$是空间速度的反对称矩阵形式。

\section{数值方法}

\subsection{数值积分}

给定初始条件$\omega_b(0)$和力矩$m_b(t)$，可以使用数值积分求解：

\textbf{欧拉法}：
\begin{align}
\dot{\omega}_b(t) &= I_b^{-1} (m_b(t) - [\omega_b(t)] I_b \omega_b(t)) \\
\omega_b(t + \Delta t) &= \omega_b(t) + \dot{\omega}_b(t) \Delta t
\end{align}

\textbf{龙格-库塔法}（更精确）：
\begin{itemize}
\item 使用四阶龙格-库塔法可以提高精度
\item 特别适用于长时间仿真
\end{itemize}

\subsection{约束处理}

如果刚体受到约束（如关节约束），需要：
\begin{itemize}
\item 添加约束力到力矩方程
\item 使用拉格朗日乘数法或投影方法
\item 确保约束条件得到满足
\end{itemize}

\section{实验验证}

\subsection{陀螺仪实验}

陀螺仪实验验证了欧拉方程中的陀螺项：
\begin{itemize}
\item 快速旋转的陀螺仪表现出进动
\item 进动角速度与角动量成正比
\item 验证了$[\omega_b] I_b \omega_b$项的存在
\end{itemize}

\subsection{自由旋转实验}

自由旋转实验验证了稳定性分析：
\begin{itemize}
\item 绕主惯性轴旋转是稳定的
\item 绕非主惯性轴旋转会产生进动
\item 角动量守恒得到验证
\end{itemize}

\section{拓展应用}

\subsection{多刚体系统}

对于多刚体系统（如机器人），每个刚体的动力学由欧拉方程描述：
\begin{equation}
m_i = I_i \dot{\omega}_i + [\omega_i] I_i \omega_i
\end{equation}

通过递归牛顿-欧拉算法可以高效计算所有关节的力矩。

\subsection{柔性体动力学}

对于柔性体，欧拉方程需要扩展：
\begin{itemize}
\item 转动惯量矩阵可能随时间变化
\item 需要考虑弹性变形的影响
\item 使用模态分析方法简化计算
\end{itemize}

\subsection{流体中的刚体}

在流体中运动的刚体，欧拉方程需要添加流体阻力项：
\begin{equation}
m_b = I_b \dot{\omega}_b + [\omega_b] I_b \omega_b + m_{\text{drag}}
\end{equation}

其中$m_{\text{drag}}$是流体阻力产生的力矩。

\section{总结}

\subsection{主要结论}

\begin{enumerate}
\item \textbf{欧拉方程}：$m_b = I_b \dot{\omega}_b + [\omega_b] I_b \omega_b$是刚体旋转动力学的基本方程。

\item \textbf{两项的物理意义}：
\begin{itemize}
\item $I_b \dot{\omega}_b$：惯性力矩，与角加速度成正比
\item $[\omega_b] I_b \omega_b$：陀螺力矩，与角速度的平方成正比
\end{itemize}

\item \textbf{主惯性轴情况}：绕主惯性轴旋转时，方程简化，陀螺项可能为零。

\item \textbf{能量关系}：功率等于旋转动能的时间导数，验证了能量守恒。

\item \textbf{稳定性}：绕主惯性轴旋转是稳定的，绕非主惯性轴旋转会产生进动。
\end{enumerate}

\subsection{关键要点}

\begin{itemize}
\item 欧拉方程是角动量定理在旋转坐标系中的体现
\item 陀螺项$[\omega_b] I_b \omega_b$是旋转坐标系效应的结果
\item 主惯性轴对齐可以简化计算和分析
\item 能量守恒是验证方程正确性的重要方法
\end{itemize}

\subsection{应用领域}

\begin{itemize}
\item 机器人动力学和控制
\item 航天器姿态控制
\item 陀螺仪和导航系统
\item 机械系统设计和分析
\item 多体系统动力学仿真
\end{itemize}

\subsection{进一步阅读}

\begin{itemize}
\item 刚体动力学经典教材
\item 机器人学教材中的动力学章节
\item 数值方法和仿真技术
\item 控制理论中的动力学建模
\end{itemize}

\end{document}